\documentclass[12pt]{article}
\usepackage[utf8]{inputenc}
\usepackage[T1]{fontenc}
\usepackage{hyperref}
\title{Tuxá Language Resources for K-12 Curriculum Development}
\author{Compiled by Claude Sonnet 5 (claude-sonnet-5) for Andr\'as Kornai}
\date{}
\begin{document}
\maketitle

\section{Introduction}
Tuxá (ISO 639-3 tud; Glottocode tuxa1239) is a dormant language isolate of Bahia, Brazil. Search results describe it as extinct as a community language since the 19th century. A 1960s research team reportedly found two women who knew about thirty words. Very little primary material exists. Current revitalization work by the Tuxá people of Rodelas focuses on Dzubukuá, a Kariri-family language the community is reclaiming, and not on Tuxá itself. A curriculum for tud would therefore have almost no native-language content and would have to rely on this related work.

\section{Primary Educational Resources}
No primers, textbooks or school materials in Tuxá turned up in the searches. The only school-related item found is a workshop page for the Capitão Francisco Rodelas Indigenous School, which concerns Dzubukuá revitalization (the page says the platform has been inactive since February 2025):\\
\url{https://idil2022-2032.org/events-activities/dzubukua-language-revitalization-project-community-based-workshop/}

\section{Linguistic and Reference Materials}
\begin{itemize}
\item Wikipedia overview of the Tuxá language: \url{https://en.wikipedia.org/wiki/Tux%C3%A1_language}
\end{itemize}
No dictionary or grammar was found.

\section{Religious and Community Resources}
\begin{itemize}
\item Durazzo (Policromias), on ritual science, the Prokár network and language revitalization in the Submédio São Francisco: \url{https://revistas.ufrj.br/index.php/policromias/article/view/46720}
\item Article in \textit{Pixo} on the effects of the flooding of old Rodelas on Tuxá relations with the \textit{encantados} and on territorial identity: \url{https://revistas.ufpel.edu.br/index.php/pixo/article/download/2867/2270}
\item Preprint found in the same search (its content was not verified): \url{https://preprints.scielo.org/index.php/scielo/preprint/view/6508}
\end{itemize}

\section{Digital and Audio Resources}
No audio, video or software resources for Tuxá were found.

\section{Curriculum Development Context}
There is no evidence of any school teaching Tuxá, and there are no speakers. The honest basis for a heritage curriculum is Tuxá history, culture and ritual life, with the community's own Dzubukuá revitalization work as the language component. Any use of Dzubukuá should be agreed with the Tuxá community and its teachers.
\end{document}